% =====================================================================
%  main.tex  -  presentation using the AIM beamer theme
%  Compile with XeLaTeX, twice (the outline is filled on the second run):
%      xelatex main.tex && xelatex main.tex
%  Files needed: beamerthemeAIM.sty and the assets/ folder.
% =====================================================================
\documentclass[aspectratio=169,11pt,t]{beamer}
\usetheme{AIM}

% -------------------------------------------------------------- metadata
\title[Forecasting Hotspots for Predictive Police Patrolling]
      {Forecasting Hotspots for\\Predictive Police Patrolling}
\author[Ismail Elabbassi]{Ismail Elabbassi}
\supervisor{Prof.\ Youness Mechqrane}
\director{Prof.\ Amal El Fallah Seghrouchni}
\date{October 2026}

\begin{document}

\titleframe
\outlineframe

% =====================================================================
\aimsection{Introduction}{Context, motivation and research questions}

\begin{frame}{Context and motivation}
  \begin{columns}[T,onlytextwidth]
    \begin{column}{0.55\textwidth}
      \begin{itemize}
        \item Placeholder point about the problem setting
        \item Why forecasting crime hotspots matters
        \begin{itemize}
          \item Sub point one
          \item Sub point two
        \end{itemize}
        \item Limits of current patrol planning
      \end{itemize}
    \end{column}
    \begin{column}{0.4\textwidth}
      \centering
      \placeholder{5.2cm}{3.6cm}{Figure placeholder\\city map, photo, \dots}
    \end{column}
  \end{columns}
\end{frame}

\begin{frame}{Research questions}
  \begin{enumerate}
    \item Can we forecast where crime will concentrate in the next period?
    \item Which features drive these forecasts?
    \item How should patrols be allocated given the forecasts?
  \end{enumerate}
  \vskip0.4cm
  \begin{block}{Main contribution}
    One sentence that states the contribution of the work.
  \end{block}
\end{frame}

% =====================================================================
\aimsection{Related Work}{Crime hotspot mapping and predictive policing}

\begin{frame}{Related work}{How this work is positioned}
  \small
  \begin{tabular}{@{}lccc@{}}
    \toprule
    \textbf{Approach} & \textbf{Spatial} & \textbf{Temporal} & \textbf{Patrol planning}\\
    \midrule
    Kernel density estimation & \checkmark & & \\
    Self-exciting point processes & \checkmark & \checkmark & \\
    Deep spatio-temporal models & \checkmark & \checkmark & \\
    \textcolor{aimBlue}{\textbf{This work}} & \checkmark & \checkmark & \checkmark\\
    \bottomrule
  \end{tabular}
\end{frame}

% =====================================================================
\aimsection{Data \& Study Area}{Crime records, features and preprocessing}

\begin{frame}{Dataset}
  \begin{columns}[T,onlytextwidth]
    \begin{column}{0.48\textwidth}
      \placeholder{6.3cm}{3.8cm}{Map of the study area}
    \end{column}
    \begin{column}{0.48\textwidth}
      \begin{itemize}
        \item Source and time span of the records
        \item Spatial grid resolution
        \item Features: history, calendar, points of interest
        \item Train / validation / test split
      \end{itemize}
    \end{column}
  \end{columns}
\end{frame}

% =====================================================================
\aimsection{Methodology}{Forecasting models and patrol allocation}

\begin{frame}{Forecasting model}
  The predicted intensity of cell $i$ at time $t$ is
  \[
    \hat{\lambda}_{i,t} = f_\theta\big(x_{i,t-1}, \dots, x_{i,t-k}\big),
  \]
  and patrols are allocated to the cells with the highest $\hat{\lambda}_{i,t}$.
  \vskip0.4cm
  \begin{exampleblock}{Intuition}
    Short explanation of the model in plain words.
  \end{exampleblock}
\end{frame}

\begin{frame}{Pipeline}
  \centering
  \placeholder{12.5cm}{4.2cm}{Pipeline diagram: data $\rightarrow$ features $\rightarrow$ model $\rightarrow$ patrol plan}
\end{frame}

% =====================================================================
\aimsection{Results}{Experiments, metrics and comparison}

\begin{frame}{Main results}
  \centering
  \metric{0.82}{Hit rate at 5\% area}\hfill
  \metric{+14\%}{Over the best baseline}\hfill
  \metric{2.3$\times$}{Predictive accuracy index}
  \vskip0.5cm
  \begin{alertblock}{Takeaway}
    One sentence that the audience should remember from the results.
  \end{alertblock}
\end{frame}

% =====================================================================
\aimsection{Conclusion}{Takeaways, limitations and future work}

\begin{frame}{Conclusion and future work}
  \begin{columns}[T,onlytextwidth]
    \begin{column}{0.48\textwidth}
      \textbf{\color{aimBlue}Takeaways}
      \begin{itemize}
        \item First takeaway
        \item Second takeaway
      \end{itemize}
    \end{column}
    \begin{column}{0.48\textwidth}
      \textbf{\color{aimBlue}Future work}
      \begin{itemize}
        \item First direction
        \item Second direction
      \end{itemize}
    \end{column}
  \end{columns}
\end{frame}

\thankyouframe{your.email@um6p.ma}

\end{document}
